% !TeX root=../main.tex
\chapter{ طراحی و ساخت حسگر چندوجهی و برد جمع‌آوری داده}
%\thispagestyle{empty} 
\section{مقدمه} 
در سال‌های اخیر، توسعه‌ی حسگرهای لمسی پیشرفته نقش مهمی در بهبود توانایی ربات‌ها برای تعامل مؤثر با محیط ایفا کرده است. در این میان، با گسترش روش های هوشمند و یادگیری عمیق، استفاده از حسگرهای چند وجهی به‌عنوان رویکردی نوین که افزونگی لازم برای روش های هوشمند را فراهم می‌کند، توجه ویژه‌ای را به خود جلب کرده است. دلیل این اهمیت آن است که هیچ نوع حسگر یکتایی قادر به پوشش کامل ویژگی‌های تماس مکانیکی نیست و ترکیب چند نوع حسگر با ماهیت متفاوت امکان استخراج اطلاعات غنی‌تر و قابل اعتمادتری را فراهم می‌سازد. برای مثال، از روش هایی که در 
\ref{sec:mechanisms}
معرفی شدند، حسگرهای فشار یا بارومتریک قادر به اندازه‌گیری نیروی استاتیکی با دقتی بالا هستند که حسگر اولیه‌ی ارائه شده در این پژوهش
\cite{sheykh2025Design}
گواه این موضوع است. اما این نوع مکانیزم بیشتر مناسب نیرو های عمودی بوده و در تفکیک نیرو و گشتاور در جهات مختلف دچار مشکل می‌شود، از طرفی در این پژوهش مشاهده شد که این حسگر ها گاهی دچار انحراف
\LTRfootnote{Drift}
 در گذر زمان می‌شوند. در مقابل، حسگر های مبتنی بر اثر هال قابلیت تفکیک مناسبی برای نیرو-گشتاور های جند محوری دارند و اغلب دچار انحراف نمی‌شوند. اما حساسیت این نوع حسگر ها به نسبت بارومتریک پایین‌تر بوده و یه شدت تحت تاثیر اغتشاشات محیطی می‌باشند، به خصوص زمانی که مواد فلزی و فرومگنتیک
 \LTRfootnote{ferromagnetic}
 در محیط حاضر باشند. از طرفی هر دو این مکانیزم ها برای نیرو ها استاتیکی مناسب بوده و پهنای باند مناسبی برای وظایفی مثل تشخیص لغزش و جنس اشیاء ندارند. برای حل چالش پهنای باند، شاید مناسب ترین مکانیزم، پیزوالکتریک باشد که مشخصاً به پدیده های دینامیکی پاسخ می‌دهد و توانایی تشخیص نیرو های استاتیکی را ندارد؛ اما استفاده از مکانیزم پیزوالکتریک به خودی خود نیازمند ادوات خاص و آزمایشگاه های مجهز است، یکی از نوآوری های دیگر این پژوهش، استفاده از میکروفون های ممز
 \LTRfootnote{MEMS}
 به عنوان نماینده‌ی مکانیزم پیزوالکتریک است. حضور این مکانیزم در مجموعه‌ی حسگر لامسه مکان ثبت ارتعاشات لحظه‌ای و اطلاعات دینامیکی تماس، مانند لغزش و بافت سطح، را فراهم می‌کند. 
 به این ترتیب، در پژوهش حاضر یک حسگر چند وجهی متشکل از یک بارومتر، دو حسگر اثر هال و یک میکروفون ممز معرفی می‌شود. ایده‌ی اصلی آن است که هر یک از این حسگرها بخشی از اطلاعات لمسی را فراهم کنند و ادغام آن‌ها منجر به بهبود دقت و افزایش قابلیت اطمینان سیستم گردد.
  
 اما حضور انواع مختلف مکانیزم با پهنای باند متفاوت، چالش جمع‌آوری داده های چندوجهی و همگم‌سازی آن‌ها را به همراه دارد. برای حل این چالش، یک سخت‌افزار اختصاصی جهت جمع‌آوری داده توسعه یافته که برای ارسال اطلاعات در لایه‌ی فیزیکی از پروتکل 
 \LR{USB}
 بهره می‌برد، این پروتکل نیاز به راحتی با عنواع سیستم ها سازگار بوده و نیاز به مبدل های سریال را از بین می‌برد. همچنین در لایه‌ی بالاتر، استفاده از کلاس
 \LR{Audio}
 و
 \LR{CDC}\LTRfootnote{Communication Device Class}
  به صورت ترکیبی امکان ارسال همزمان داده‌های صوتی و اطلاعات دیگر حسگر ها را فراهم می‌سازد.

با توجه به اهمیت چندوجهی بودن داده‌ها، بخشی از این فصل به فرایند کالیبراسیون اختصاص داده شده است. در این فرایند، ابتدا کالیبراسیون استاتیک جهت همسان‌سازی پاسخ اولیه‌ی حسگر ها انجام می‌شود و سپس مدل‌های دینامیکی برای توصیف رفتار غیرخطی و وابسته به زمان حسگر ارائه خواهند شد. در نهایت، نتایج کالیبراسیون و مدل‌سازی به‌همراه نمودارها و تحلیل‌ها ارائه می‌شود تا نشان داده شود که چگونه ادغام چند نوع مکانیزم می‌تواند به طراحی یک حسگر لمسی دقیق و کارا منجر شود.
\section{قطعات و مکانیزم های استفاده شده}
در ابن بخش به معرفی قطعات به کار رفته در حسگر لامسه می‌پردازیم. یکی از معیار ها و محدودیت های اصلی در انتخاب قطعات، موجود بودن آن‌ها در بازار ایران هنگام طراحی و ساخت حسگر پیشنهاد شده یوده است. بنابر این، در پژوهش های آتی با توجه وجود قطعات مناسب تر در بازار امکان ساخت حسگر لامسه با عملکری بهتر ممکن خواهد بود.
\begin{figure}[ht]
	\centering 
	\subfloat[حسگر فشار \LR{MPL3115A2}]{ 
		\includegraphics[width=0.35\textwidth]{MPL3115A2}}
	%\hspace{2mm}
	\subfloat[حسگر اثر هال \LR{MMC5883MA} و \LR{QMC5883L}]{ 
		\includegraphics[width=0.35\textwidth]{MMC5883}}%
	
	\subfloat[میکروفون ممز \LR{ICS43434}]{ 
		\includegraphics[width=0.4\textwidth]{isc43434}}%
	\caption{حسگر های استفاده شده.}
	\label{fig:parts1} %% label for entire figure
\end{figure}
\\
\subsection{حسگر بارومتریک \LR{MPL3115A2}}
	
	 
حسگر 
\LR{MPL3115A2}
 یک حسگر فشار مطلق دیجیتال از شرکت \LR{NXP} است که قادر به اندازه‌گیری فشار در بازه‌ی 0 تا 125 کیلوپاسکال با دقت 1 پاسکال می‌باشد. هرچند که در آزمایش های عملی، مقادیر نزدیک به 140 کیلوپاسکال در خروجی حسگر نیز مشاهده شد. این دقت بالا موجب شده است که بتوان تغییرات بسیار کوچک نیروهای وارد بر سطح حسگر را آشکارسازی کرد. خروجی حسگر به‌صورت دیجیتال و از طریق پروتکل \LR{I2C} در دسترس است. ایده‌ی استفاده از این حسگر در طراحی حاضر دقت بالا در اندازه‌گیر نیروی عمودی بوده است و به دلیل حساسیت بالای خود نسبت به تغییرات فشار، امکان تشخیص تغییرات دقیق نیروی وارد بر لایه‌ی سیلیکونی را فراهم می‌کند. محدودیت اصلی این حسگر نیاز به اصلاحات مکانیکی و احتمال وجود انحراف در گذر زمان است که در مراحل بعدی با کالیبراسیون جبران خواهد شد. زمان نمونه برداری برای این حسگر حداقل 6 میلی‌ثانیه می‌باشد که برای کاربرد های استاتیک و حتی شبه استاتیک مناسب است. 
\subsection{حسگر اثر هال \LR{MMC5883MA}}


حسگر 
\LR{MMC5883MA}
 یک مغناطیس‌سنج سه‌محوره دقیق از شرکت \LR{MEMSIC} است. این حسگر توانایی اندازه‌گیری میدان مغناطیسی در بازه‌ی
  $\pm 8$ 
  گوس با وضوح حدود 0.25 میلی‌گاوس را دارد. پروتکل ارتباطی \LR{I2C} بوده و زمان نمونه برداری می‌تواند تا حداقل 1 میلی‌ثانیه انتخاب شود.همانطور که در شکل 
 \ref{fig:components3d}
  نشان داده شده، این حسگر در زیر یکی از قطب‌های آهنربای رینگی شعاعی نصب شده است، به‌طوری ‌که خطوط میدان آهنربا تقریباً عمود بر حسگر باشد. این جاگذاری باعث می‌شود که تغییرات میدان در محور Z حسگر به دلیل جابجایی اندک آهنربا در اثر نیروی مماسی  بیشینه شود. بنابراین، \LR{MMC5883MA} حساسیت بیشتری نسیت به نیرو های مماسی داشته و نقش کلیدی در ثبت مولفه‌های برشی نیرو ایفا می‌کند. 
 \subsection{حسگر اثر هال \LR{QMC5883L}}
 
 حسگر \LR{QMC5883L} یک مغناطیس‌سنج سه‌محوره با دقت متوسط از شرکت \LR{QST} است.بازه‌ی کاری این حسگر در تا 
  $\pm 8$ 
  گوس می‌تواند تنظیم شود و دقت آن حدود 1 تا 2 میلی‌گاوس است. پروتکل ارتباطی آن \LR{I2C} می‌باشد و حداقل زمان نمونه‌برداری 5 میلی‌ثانیه را داراست. در طراحی حاضر،این حسگر در مرکز آهنربای رینگی و در راستای محور عمودی قرار داده شده است. این موقعیت باعث می‌شود حسگر به تغییرات میدان مغناطیسی ناشی از چرخش یا گشتاور حول محور عمودی حساسیت بیشتری داشته باشد. بنابراین \LR{QMC5883L} برای آشکارسازی لنگر انداختن جسم یا گشتاورهای پیچشی بسیار مناسب است و مکمل \LR{MMC5883MA} محسوب می‌شود. دلیل انتخاب مدل های متفاوت برای حسگر اثر هال، صرفاً، تفاوت در آدرس دهی و قابلیت تبادل اطلاعات در یک خط داده بوده است.
 
 \subsection{میکروفون  \LR{ISC43434}}
 میکروفون \LR{ICS43434} یک میکروفون دیجیتال ممز با خروجی 
 \LR{I2S}\LTRfootnote{Inter-Integrated Circuit Sound } 
 از شرکت \LR{TDK} بوده و قابلیت ضبط صوتی با نرخ نمونه‌برداری تا 48 کیلوهرتز و وضوح 24 بیت را داراست. این میکروفون دارای نویز پایین 
 $65 \: dBA$
  و حساسیت حدود 
  $-26 dB FS$
  بوده و در بازه‌ی فرکانسی 60 هرتز تا 20 کیلوهرتز کار می‌کند. در این پروژه، میکروفون در زیر لایه‌ی کشسان سیلیکونی و در کنار سایر حسگرها قرار گرفته و نقش آن اندازه‌گیری ارتعاشات دینامیکی، لغزش و شناسایی بافت سطح است. استفاده از رابط \LR{I2S} و بهره‌گیری از 
  \LR{DMA}\LTRfootnote{Direct Memory Access}
   امکان خواندن پیوسته داده‌ها را بدون بارگذاری زیاد روی پردازنده فراهم می‌سازد.
\begin{figure}[t]
	\centering
	\centerline{\includegraphics[width=0.7\textwidth]{components_placement}}
	\caption{
		جاگذاری حسگر ها و آهنربا. یک حسگر فشار و اثر هال روی برد و یک حسگر اثر هال و میکروفون در زیر برد قرار دارند.
	}
	\label{fig:components3d}
\end{figure} 

\subsection{آهنربا}
به منظور استخراج ویژگی‌های مربوط به تغییرات میدان مغناطیسی ناشی از جایجایی یا تغییر شکل ناحیه‌ی کشسان و تبدیل آن به نیرو، از یک آهنربای حلقوی نئودیمیومی با قطر داخلی ۵ میلی‌متر، قطر خارجی ۱۲ میلی‌متر و ضخامت ۳ میلی‌متر استفاده شده است. این آهنربا با گرید N35 تولید شده و به دلیل چگالی شار مغناطیسی بالا و پایداری در ابعاد کوچک، گزینه‌ای مناسب برای ترکیب با حسگرهای هال محسوب می‌شود. در این سامانه، تغییر موقعیت نسبی آهنربا نسبت به حسگرهای هال، تغییراتی در میدان مغناطیسی ثبت ‌شده ایجاد کرده که امکان استنتاج  نیروهای اعمالی را فراهم می‌سازد.